\documentclass[11pt]{article}
\usepackage[margin=1in]{geometry}
\usepackage{amsmath,amssymb,amsthm}
\usepackage{booktabs}
\usepackage{hyperref}

\newtheorem{theorem}{Theorem}
\newtheorem{definition}{Definition}
\newtheorem{remark}{Remark}

\title{A Point-Transitive Kirkman Triple System of Order 9:\\
Disproof of Conjecture 00000008420 (Transitive-KTS Clause)}
\author{TLMC submission}
\date{September 2026}

\begin{document}
\maketitle

\begin{abstract}
Conjecture 00000008420 asserts, among other clauses, that ``the smallest order
of a KTS with a transitive automorphism is the point-transitive type of order
$15$''.  We exhibit an explicit counterexample: the line system of the affine
plane $\mathrm{AG}(2,3)$ is a resolvable Steiner triple system of order $9$
(a $\mathrm{KTS}(9)$) whose translation group $\mathbb{F}_3^2$ acts
transitively on the points.  Since $9 < 15$, the clause is false.
The construction is verified by independent enumeration
(\texttt{reproduce.py}) and by an axiom-free Lean~4 formalization in which all
statements are proved by \texttt{decide} over explicit finite enumerations.
\end{abstract}

\section{Scope of the disproof}

The conjecture under attack is a conjunction of four clauses.  We refute only
the third one:
\begin{center}
\emph{``the smallest order of a KTS with a transitive automorphism is the
point-transitive type of order $15$''}.
\end{center}
The other clauses are consistent with our counterexample and no claim is made
about them: the order $9 \equiv 3 \pmod 6$ (clause~1), and the number of
parallel classes of the system below is exactly $4 = (9-1)/2$ (clause~2).
Clause~4 (no KTS with cyclic automorphism group of prime order) is untouched.

\section{Definitions}

\begin{definition}[Steiner triple system]
A \emph{Steiner triple system} $\mathrm{STS}(n)$ is a pair $(P,\mathcal{L})$
where $P$ is a set of $n$ points and $\mathcal{L}$ is a collection of
$3$-element subsets of $P$ (blocks) such that every pair of distinct points
lies in exactly one block.
\end{definition}

\begin{definition}[Parallel class; resolvability]
Let $(P,\mathcal{L})$ be an STS.  A \emph{parallel class} is a subset
$\mathcal{C} \subseteq \mathcal{L}$ of pairwise disjoint blocks whose union is
$P$.  The system is \emph{resolvable} if $\mathcal{L}$ is partitioned into
parallel classes.  A \emph{Kirkman triple system} $\mathrm{KTS}(n)$ is a
resolvable $\mathrm{STS}(n)$.
\end{definition}

\begin{definition}[Automorphism; point-transitivity]
An \emph{automorphism} of $(P,\mathcal{L})$ is a bijection
$\alpha : P \to P$ with $\alpha(\mathcal{L}) = \mathcal{L}$, i.e.\
$\{ \alpha(B) : B \in \mathcal{L} \} = \mathcal{L}$.  The system is
\emph{point-transitive} if its automorphism group acts transitively on $P$.
\end{definition}

\section{The construction: $\mathrm{KTS}(9)$ as the lines of $\mathrm{AG}(2,3)$}

Let $\mathbb{F}_3 = \{0,1,2\}$ and let the point set be
$P = \mathbb{F}_3^2$, so $|P| = 9$.  Define $12$ lines:
\begin{itemize}
\item for each slope $m \in \mathbb{F}_3$ and intercept $b \in \mathbb{F}_3$,
\[
L_{m,b} \;=\; \{\, (x,\; mx+b) : x \in \mathbb{F}_3 \,\},
\]
giving $9$ non-vertical lines;
\item for each $c \in \mathbb{F}_3$, the vertical line
\[
V_c \;=\; \{\, (c, y) : y \in \mathbb{F}_3 \,\},
\]
giving $3$ further lines.
\end{itemize}
Table~\ref{tab:lines} lists all $12$ lines explicitly; the index $k$ used
there ($k = 3m+b$ for the slope lines, $k = 9+c$ for the vertical lines) is
the index used by the Lean formalization and the script.

\begin{table}[ht]
\centering
\begin{tabular}{cccc}
\toprule
$k$ & type & points & parallel class \\
\midrule
0 & $m{=}0,\, b{=}0$ & $(0,0)\ (1,0)\ (2,0)$ & $C_0$ \\
1 & $m{=}0,\, b{=}1$ & $(0,1)\ (1,1)\ (2,1)$ & $C_0$ \\
2 & $m{=}0,\, b{=}2$ & $(0,2)\ (1,2)\ (2,2)$ & $C_0$ \\
3 & $m{=}1,\, b{=}0$ & $(0,0)\ (1,1)\ (2,2)$ & $C_1$ \\
4 & $m{=}1,\, b{=}1$ & $(0,1)\ (1,2)\ (2,0)$ & $C_1$ \\
5 & $m{=}1,\, b{=}2$ & $(0,2)\ (1,0)\ (2,1)$ & $C_1$ \\
6 & $m{=}2,\, b{=}0$ & $(0,0)\ (1,2)\ (2,1)$ & $C_2$ \\
7 & $m{=}2,\, b{=}1$ & $(0,1)\ (1,0)\ (2,2)$ & $C_2$ \\
8 & $m{=}2,\, b{=}2$ & $(0,2)\ (1,1)\ (2,0)$ & $C_2$ \\
9 & vertical, $c{=}0$ & $(0,0)\ (0,1)\ (0,2)$ & $C_3$ \\
10 & vertical, $c{=}1$ & $(1,0)\ (1,1)\ (1,2)$ & $C_3$ \\
11 & vertical, $c{=}2$ & $(2,0)\ (2,1)\ (2,2)$ & $C_3$ \\
\bottomrule
\end{tabular}
\caption{The $12$ lines of the system and the partition into parallel classes.}
\label{tab:lines}
\end{table}

\section{Verification of the properties}

Each property below was checked twice, independently: (i) by a
dependency-free Python enumeration (\texttt{reproduce.py}) that builds the 12
lines explicitly and enumerates all $36$ point-pairs, the $4$ parallel
classes, the $9$ translations and all pairs of points; and (ii) by an
axiom-free Lean~4 formalization (\texttt{lean4/Main.lean}, toolchain
\texttt{leanprover/lean4:v4.33.1}) whose theorems \texttt{sts\_property},
\texttt{resolvable}, \texttt{translation\_invariant}, \texttt{transitive} and
\texttt{refute} are proved by \texttt{decide} over the explicit finite
enumerations and depend on no axioms (no \texttt{sorryAx}, no
\texttt{Classical.choice}, not even \texttt{propext}).

\subsection{Steiner triple system}

There are $\binom{9}{2} = 36$ unordered pairs of distinct points.  The 12
lines contain $12 \cdot \binom{3}{2} = 36$ point-pairs in total, and the
enumeration of Table~\ref{tab:lines} shows that no two lines share more than
one point, i.e.\ all $36$ line-pairs are distinct.  Hence every pair of
distinct points lies in exactly one line: $(P, \mathcal{L})$ is an
$\mathrm{STS}(9)$.  (Equivalently: for each of the $81$ ordered pairs
$(p,q)$ with $p \neq q$, exactly one line contains both; each point lies on
exactly $4$ lines.)

\subsection{Resolvability}

The four parallel classes are
\[
C_0 = \{0,1,2\}, \qquad
C_1 = \{3,4,5\}, \qquad
C_2 = \{6,7,8\}, \qquad
C_3 = \{9,10,11\},
\]
using the line indices of Table~\ref{tab:lines}.  Within each class the three
lines are pairwise disjoint and their union is all of
$\mathbb{F}_3^2$: this is the elementary fact that over $\mathbb{F}_3$ the
lines of fixed slope $m$ (respectively the vertical lines) partition the
plane.  The four classes are disjoint and cover all $12$ lines, so the system
is resolvable: it is a $\mathrm{KTS}(9)$.  The number of parallel classes is
$4 = (9-1)/2$, matching clause~2 of the conjecture.

\subsection{Translations are automorphisms}

For $v \in \mathbb{F}_3^2$ let $t_v(p) = p + v$ (addition in
$\mathbb{F}_3^2$).  If $p$ lies on $L_{m,b}$ then $t_v(p)$ lies on
$L_{m,\, b + m v_1 + v_2}$; if $p$ lies on $V_c$ then $t_v(p)$ lies on
$V_{c + v_1}$.  Thus the image of a line under $t_v$ is again a line; since
$t_v$ is a bijection of $P$, the map induced on the line set is a bijection
of $\mathcal{L}$, i.e.\ $t_v$ is an automorphism.  This was verified by
explicit enumeration for all $9$ translations and all $12$ lines (checking
for each $(v,k)$ that a line $k'$ with
$p \in L_k \Leftrightarrow t_v(p) \in L_{k'}$ for all $p$ exists).

\subsection{Point-transitivity}

For all $p, q \in \mathbb{F}_3^2$ the translation $t_{q-p}$ maps $p$ to $q$:
\[
t_{q-p}(p) = p + (q - p) = q.
\]
Hence the automorphism group contains the translation subgroup
$\{t_v : v \in \mathbb{F}_3^2\} \cong \mathbb{F}_3^2$ of order $9$, which acts
transitively on the $9$ points.  The system is point-transitive.

\section{Conclusion}

\begin{theorem}
There exists a point-transitive Kirkman triple system of order $9$, namely
the line system of $\mathrm{AG}(2,3)$.
\end{theorem}

\begin{theorem}[Disproof of clause 3 of conjecture 00000008420]
The smallest order of a KTS with a transitive automorphism is at most $9$,
and $9 < 15$.  Therefore the clause ``the smallest order of a KTS with a
transitive automorphism is the point-transitive type of order $15$'' is
false.
\end{theorem}

\begin{remark}
The full automorphism group of $\mathrm{KTS}(9)$ as constructed is
$\mathrm{AGL}(2,3)$, of order $9 \cdot 48 = 432$; only the transitivity of
its translation subgroup is used above.
\end{remark}

\end{document}
